\chapter{Cómo usar este libro}

\section*{¿Para quién es?}

Para cualquier persona que quiera entender, desde cero, cómo funciona un dron que busca personas por
sí solo: estudiantes, ingenieros de otras áreas o curiosos. No necesitas saber de drones, de visión
computacional ni de redes neuronales. Ayuda saber usar una computadora con Linux y recordar algo de
física y matemáticas de bachillerato; cuando haga falta algo más, se explica en el momento.

El libro cuenta una historia concreta: una simulación en la que un dron Tarot 650 despega, recorre en
franjas un bosque de $120\times120$ metros, mira hacia abajo con su cámara, reconoce a las personas
con una red neuronal (YOLO), calcula en qué punto del suelo está cada una, avisa sus coordenadas GPS y
aterriza. El vuelo lo controla un piloto automático real (PX4 o ArduPilot) o una red neuronal
entrenada para imitarlo. Cada pieza de esa historia es una excusa para explicar un concepto.

\section*{Cómo está organizado}

\begin{description}
  \item[Parte I. Lo básico.] Qué es un dron y cómo vuela, la física mínima, qué es una simulación,
        y ROS~2 y PX4, que conectan todas las piezas.
  \item[Parte II. Ver y ubicar.] Qué es una imagen, cómo aprende una red neuronal a reconocer
        personas y cómo se pasa de un recuadro en la imagen a un punto en el mapa, con su incertidumbre.
  \item[Parte III. El control.] Cómo se mantiene el dron en su ruta: control en cascada y una red
        neuronal que aprende a volar imitando al piloto automático.
  \item[Parte IV. Construir la simulación.] La instalación y cómo se metieron en Gazebo el Tarot 650, su
        cámara, el bosque y las personas.
  \item[Parte V. La misión.] Cómo el dron decide qué hacer: el barrido, la re-observación de casos
        dudosos, los avisos GPS y la evasión de obstáculos, con PX4 o con ArduPilot.
  \item[Parte VI. Experimentos y manos a la obra.] Cómo se diseña y analiza una campaña de cientos de
        vuelos, un tutorial para tu primer vuelo y la referencia de órdenes.
\end{description}

Al final hay un \textbf{glosario}. Si solo quieres usar la simulación, ve directo al
capítulo~\ref{cap:primer-vuelo}.

\section*{Los recuadros}

\begin{recuerda}
Lo más importante del tema, en pocas palabras.
\end{recuerda}
\begin{consejo}
Trucos prácticos que ahorran tiempo o problemas.
\end{consejo}
\begin{ejemplo}
Una comparación con algo cotidiano.
\end{ejemplo}
\begin{tecnico}
Detalles matemáticos o de programación que puedes saltarte la primera vez.
\end{tecnico}
\begin{aviso}
Errores fáciles de cometer.
\end{aviso}
\begin{leccion}
Problemas reales que aparecieron durante el desarrollo y cómo se resolvieron.
\end{leccion}
\begin{ejercicio}
Algo que puedes hacer en la simulación para comprobar lo que acabas de leer.
\end{ejercicio}

\section*{Convenciones}

\begin{itemize}
  \item Las órdenes de terminal van en recuadros grises:
\begin{lstlisting}[style=bash]
ros2 launch mision simulacion.launch.py
\end{lstlisting}
  \item Nombres de archivos, programas y tópicos en letra de máquina: \cod{localizador_node}.
  \item \cod{~} es tu carpeta personal. El paquete principal está en \cod{~/ros2_ws/src/mision}.
  \item Distancias en metros, tiempos en segundos, ángulos en radianes o grados (se indica).
\end{itemize}
